% VAL-019 -- generado por herramientas/gen_casos.py; editar el texto libremente.
\section{VAL-019: P-Delta: benchmark AISC Caso 2 (voladizo con $H$)}\label{sec:VAL-019}
\begin{ficha}
Fuente & \cite{aisc16}, Comentario C2, Fig.~C-C2.3, p.~16.1-291; solución cerrada propia \cite{timoshenko}\\
Ticket & MAM-104\\
Engine & commit \texttt{\EngineCommit}\\
Binario & SHA-256 \texttt{\EngineBinSHA}\\
Entrada & SHA-256 \texttt{\EntradaSHA{VAL-019}}\\
Tolerancia & $|c-r|\le A+R|r|$, $A=0$, $R=0.03$ (clase consolidada; orden de ejecución DESCONOCIDO; el marco también distingue preregistro, DESCONOCIDO y VAL-037 tras la r1 de \#23)\\
Estado de datos & \directlua{mamut.estado('VAL-019')}\\
\end{ficha}

\paragraph{Modelo.}
\begin{center}\diag{voladizoP}\end{center}

\paragraph{Referencia y derivación.} Voladizo W14$\times$48 de $L=336$~in con carga lateral $H=1$~kip y compresión $P$ en la punta (geometría de \cite{aisc16}, Fig.~C-C2.3). Con $k=\sqrt{P/EI}$ (derivación propia, efectos $P$-$\Delta$ y $P$-$\delta$, sin cortante):
\[
M_\text{base}=\frac{H\tan kL}{k},\qquad \Delta_\text{tip}=\frac{H\,(\tan kL-kL)}{P\,k},
\]
con los límites $HL$ y $HL^3/3EI$ cuando $P\to0$. Valores: \texttt{herramientas/gen\_referencia.py}.


\paragraph{Reproducción.}
\reproducir{VAL-019}

% Q2-PLACEHOLDER INICIO (VAL-019)
% Solo Q2 escribe datos/VAL-019/opensees.csv y el texto de modelo OpenSees de este bloque,
% en su propio PR. J2 no lee ni escribe valores de OpenSees.
\paragraph{Modelo de contraste (OpenSees).} \QdosPendiente
% Q2-PLACEHOLDER FIN (VAL-019)

\paragraph{Resultados.}
\begin{center}
\begin{tikzpicture}
\begin{axis}[grafico, xlabel={$P$ [kN]}, ylabel={$M_\text{base}$ [kN\,m]}]
\directlua{mamut.serie('VAL-019','referencia','P_kN','M_kNm','ref','Referencia analítica')}
\directlua{mamut.serie('VAL-019','engine','P_kN','M_kNm','eng','MAMUT Engine')}
\directlua{mamut.serie('VAL-019','opensees','P_kN','M_kNm','ops','OpenSees (contraste)')}
\end{axis}
\end{tikzpicture}
\end{center}
\begin{center}\small
\begin{tabular}{rrrrrrl}\toprule
$P$ [kN] & Ref. & Engine & OpenSees & Error rel. & Tol. & Veredicto\\\midrule
\directlua{mamut.tabla('VAL-019','P_kN','M_kNm',0,0.03)}
\bottomrule\end{tabular}
\end{center}
